\documentclass[12pt,a4paper]{article}

% -------------------------------------------------
% PACKAGES
% -------------------------------------------------
\usepackage[margin=1in]{geometry}
\usepackage{amsmath,amssymb}
\usepackage{booktabs}
\usepackage{array}
\usepackage{graphicx}
\usepackage{float}
\usepackage{xcolor}
\usepackage{listings}
\usepackage{hyperref}
\usepackage{fancyhdr}
\usepackage{setspace}
\usepackage{titlesec}

% -------------------------------------------------
% DOCUMENT SETTINGS
% -------------------------------------------------
\onehalfspacing

\hypersetup{
    colorlinks=true,
    linkcolor=blue,
    urlcolor=blue,
    citecolor=blue
}

\pagestyle{fancy}
\fancyhf{}
\rhead{Chapter 4: Classification}
\lhead{Statistical Learning}
\cfoot{\thepage}

\titleformat{\section}
{\Large\bfseries}
{\thesection}
{1em}
{}

\titleformat{\subsection}
{\large\bfseries}
{\thesubsection}
{1em}
{}

% -------------------------------------------------
% PYTHON CODE STYLE
% -------------------------------------------------
\definecolor{codegreen}{rgb}{0,0.5,0}
\definecolor{codegray}{rgb}{0.45,0.45,0.45}
\definecolor{codepurple}{rgb}{0.45,0,0.7}

\lstset{
    language=Python,
    basicstyle=\ttfamily\small,
    keywordstyle=\color{blue},
    commentstyle=\color{codegreen},
    stringstyle=\color{codepurple},
    numberstyle=\tiny\color{codegray},
    numbers=left,
    numbersep=8pt,
    frame=single,
    breaklines=true,
    showstringspaces=false,
    tabsize=4
}

% -------------------------------------------------
% TITLE PAGE
% -------------------------------------------------

\begin{document}

\begin{titlepage}

\centering

\vspace*{1.5cm}

{\Large \textbf{STATISTICAL LEARNING}}\\[0.5cm]

{\Huge \textbf{Chapter 4: Classification}}\\[0.5cm]

{\Large Exercises 5, 10, 14 and 16}\\[1.5cm]

\rule{\textwidth}{0.5pt}\\[0.5cm]

{\Large \textbf{An Introduction to Statistical Learning\\
with Applications in Python}}\\[0.5cm]

\rule{\textwidth}{0.5pt}

\vfill

\begin{tabular}{rl}
\textbf{Name:} & Eliane UMWALI \\[0.3cm]
\textbf{Student ID:} & 101206 \\[0.3cm]
\textbf{Chapter:} & 4 -- Classification \\[0.3cm]
\textbf{Software:} & Python / Jupyter Notebook \\[0.3cm]
\textbf{Book:} & An Introduction to Statistical Learning with Applications in Python
\end{tabular}

\vfill

{\large \today}

\end{titlepage}

% -------------------------------------------------
% CONTENTS
% -------------------------------------------------

\tableofcontents
\newpage

% =================================================
% INTRODUCTION
% =================================================

\section{Introduction}

Classification is a supervised statistical learning problem in which the
response variable is qualitative rather than quantitative. The objective is
to use one or more predictors to assign an observation to one of two or
more classes.

Chapter 4 of \textit{An Introduction to Statistical Learning with Applications
in Python} introduces several important classification methods, including
logistic regression, linear discriminant analysis (LDA), quadratic
discriminant analysis (QDA), naive Bayes, and K-nearest neighbors (KNN).

This report addresses Exercises 5, 10, 14, and 16 from the chapter. The
first two exercises focus on the theoretical differences between LDA and
QDA and the derivation of the LDA discriminant function for one predictor.
Exercises 14 and 16 apply classification methods to the Auto and Boston
datasets.

For the computational exercises, the analysis was performed in Python using
the ISLP datasets and scikit-learn classification methods.

% =================================================
% EXERCISE 5
% =================================================

\section{Exercise 5: LDA and QDA}

\subsection{(a) Linear Bayes Decision Boundary}

If the Bayes decision boundary is linear, QDA is expected to perform better
on the training set because QDA is more flexible than LDA. Its greater
flexibility allows it to fit the training observations at least as well as
the more restricted LDA model.

However, on the test set, LDA is expected to perform better. When the true
Bayes decision boundary is linear, the additional flexibility of QDA is not
necessary. QDA estimates a separate covariance matrix for each class, which
results in more parameters and higher variance.

Therefore:

\begin{itemize}
    \item \textbf{Training set:} QDA is expected to perform better.
    \item \textbf{Test set:} LDA is expected to perform better.
\end{itemize}

The reason is the bias--variance trade-off. LDA is less flexible and
therefore has lower variance, while QDA has greater flexibility but higher
variance.

\subsection{(b) Non-linear Bayes Decision Boundary}

If the Bayes decision boundary is non-linear, QDA is expected to perform
better on both the training and test sets.

LDA assumes a common covariance matrix across the classes and consequently
produces a linear decision boundary. QDA allows each class to have its own
covariance matrix and therefore produces a quadratic decision boundary.

Thus, when the true boundary is non-linear, QDA can represent the underlying
relationship more accurately.

\begin{itemize}
    \item \textbf{Training set:} QDA is expected to perform better.
    \item \textbf{Test set:} QDA is expected to perform better, provided
    there are enough observations to estimate the additional parameters
    reliably.
\end{itemize}

\subsection{(c) Effect of Increasing Sample Size}

As the sample size $n$ increases, we expect the test prediction accuracy of
QDA relative to LDA to improve.

QDA estimates a separate covariance matrix for each class and consequently
estimates more parameters than LDA. With a small sample size, these estimates
can have high variance. As the sample size increases, the parameter
estimates become more stable and the variance of QDA decreases.

Therefore, QDA becomes more attractive relative to LDA as the sample size
increases.

\subsection{(d) True or False}

\textbf{False.}

Although QDA is flexible enough to model a linear decision boundary, this
additional flexibility does not guarantee a lower test error.

When the true Bayes decision boundary is linear, LDA already provides the
appropriate form of decision boundary. QDA estimates more parameters and
therefore has higher variance. The additional flexibility can lead to
overfitting without providing a corresponding reduction in bias.

Therefore, when the true boundary is linear, LDA will generally have better
test performance than QDA.

% =================================================
% EXERCISE 10
% =================================================

\section{Exercise 10: LDA When $p=1$}

For LDA, the log posterior odds between class $k$ and the reference class
$K$ can be written as

\[
\log
\left(
\frac{\Pr(Y=k\mid X=x)}
{\Pr(Y=K\mid X=x)}
\right)
=
a_k+\sum_{j=1}^{p}b_{kj}x_j.
\]

When $p=1$, the mean vectors become scalars $\mu_k$ and $\mu_K$, while the
shared covariance matrix becomes the scalar variance $\sigma^2$.

The expression becomes

\[
\log
\left(
\frac{\Pr(Y=k\mid X=x)}
{\Pr(Y=K\mid X=x)}
\right)
=
\log\left(\frac{\pi_k}{\pi_K}\right)
-\frac{(x-\mu_k)^2}{2\sigma^2}
+\frac{(x-\mu_K)^2}{2\sigma^2}.
\]

Expanding the squared terms gives

\[
\begin{aligned}
&\log\left(\frac{\pi_k}{\pi_K}\right)
-\frac{x^2-2x\mu_k+\mu_k^2}{2\sigma^2}
+\frac{x^2-2x\mu_K+\mu_K^2}{2\sigma^2}.
\end{aligned}
\]

The $x^2$ terms cancel, giving

\[
\log\left(\frac{\pi_k}{\pi_K}\right)
-\frac{\mu_k^2-\mu_K^2}{2\sigma^2}
+
\frac{\mu_k-\mu_K}{\sigma^2}x.
\]

Comparing this with

\[
a_k+b_{k1}x,
\]

we obtain

\[
\boxed{
a_k =
\log\left(\frac{\pi_k}{\pi_K}\right)
-\frac{\mu_k^2-\mu_K^2}{2\sigma^2}
}
\]

and

\[
\boxed{
b_{k1} =
\frac{\mu_k-\mu_K}{\sigma^2}
}.
\]

Thus, when there is only one predictor, the LDA log posterior odds are a
linear function of $x$.

% =================================================
% EXERCISE 14
% =================================================

\section{Exercise 14: Auto Dataset}

The objective of this exercise is to predict whether a vehicle has high or
low gas mileage using the Auto dataset.

\subsection{(a) Creating \texttt{mpg01}}

The first step is to create a binary response variable called
\texttt{mpg01}. The median of MPG is used as the threshold.

Vehicles with MPG above the median are assigned a value of 1, while vehicles
with MPG below the median are assigned a value of 0.

The Python implementation was:

\begin{lstlisting}
from ISLP import load_data

Auto = load_data('Auto')

median_mpg = Auto['mpg'].median()

Auto['mpg01'] = (
    Auto['mpg'] > median_mpg
).astype(int)

print("Median MPG:", median_mpg)
\end{lstlisting}

The median MPG obtained from the dataset was

\[
\boxed{22.75}.
\]

Therefore,

\[
\texttt{mpg01}=1
\]

represents vehicles with MPG above 22.75, while

\[
\texttt{mpg01}=0
\]

represents vehicles with MPG below or equal to the median.

\subsection{(b) Graphical Exploration and Predictor Selection}

The relationship between \texttt{mpg01} and the available predictors was
examined using graphical and numerical summaries.

The strongest associations were observed for:

\begin{itemize}
    \item cylinders,
    \item weight,
    \item displacement,
    \item horsepower,
    \item year, and
    \item origin.
\end{itemize}

The strongest negative associations were observed for cylinders, weight,
displacement, and horsepower. This indicates that vehicles with larger
engines, more cylinders, greater horsepower, and greater weight tend to
have lower fuel efficiency.

Year showed a positive association with \texttt{mpg01}, suggesting that
newer vehicles tend to have better fuel efficiency.

Based on these findings, the following predictors were used in the
classification models:

\[
\boxed{
\text{cylinders, displacement, horsepower, weight, year}
}
\]

\subsection{(c) Training and Test Sets}

The data were divided into training and test sets using a 50--50 split.
The random state was set to 1 so that the results could be reproduced.

\begin{lstlisting}
from sklearn.model_selection import train_test_split

features = [
    'cylinders',
    'displacement',
    'horsepower',
    'weight',
    'year'
]

X = Auto[features]
y = Auto['mpg01']

X_train, X_test, y_train, y_test = train_test_split(
    X,
    y,
    test_size=0.5,
    random_state=1,
    stratify=y
)
\end{lstlisting}

The resulting datasets contained 196 observations each.

\subsection{(d) Linear Discriminant Analysis}

LDA was fitted using the selected predictors.

\begin{lstlisting}
from sklearn.discriminant_analysis import LinearDiscriminantAnalysis
from sklearn.metrics import accuracy_score

lda = LinearDiscriminantAnalysis()

lda.fit(X_train, y_train)

lda_pred = lda.predict(X_test)

lda_error = 1 - accuracy_score(y_test, lda_pred)

print("LDA test error:", lda_error)
\end{lstlisting}

The test error obtained was approximately

\[
\boxed{9.69\%}.
\]

Therefore, the corresponding test accuracy was approximately

\[
\boxed{90.31\%}.
\]

\subsection{(e) Quadratic Discriminant Analysis}

QDA was fitted using the same predictors.

\begin{lstlisting}
from sklearn.discriminant_analysis import QuadraticDiscriminantAnalysis

qda = QuadraticDiscriminantAnalysis()

qda.fit(X_train, y_train)

qda_pred = qda.predict(X_test)

qda_error = 1 - accuracy_score(y_test, qda_pred)

print("QDA test error:", qda_error)
\end{lstlisting}

The test error was approximately

\[
\boxed{11.73\%}.
\]

The corresponding test accuracy was approximately

\[
\boxed{88.27\%}.
\]

QDA performed slightly worse than LDA on this particular split, suggesting
that the additional flexibility of QDA was not beneficial for this dataset.

\subsection{(f) Logistic Regression}

Logistic regression was fitted as follows:

\begin{lstlisting}
from sklearn.linear_model import LogisticRegression

logit = LogisticRegression(max_iter=5000)

logit.fit(X_train, y_train)

logit_pred = logit.predict(X_test)

logit_error = 1 - accuracy_score(y_test, logit_pred)

print("Logistic regression test error:", logit_error)
\end{lstlisting}

The test error was approximately

\[
\boxed{10.20\%},
\]

giving a test accuracy of approximately

\[
\boxed{89.80\%}.
\]

\subsection{(g) Naive Bayes}

A Gaussian naive Bayes classifier was fitted using the selected predictors.

\begin{lstlisting}
from sklearn.naive_bayes import GaussianNB

nb = GaussianNB()

nb.fit(X_train, y_train)

nb_pred = nb.predict(X_test)

nb_error = 1 - accuracy_score(y_test, nb_pred)

print("Naive Bayes test error:", nb_error)
\end{lstlisting}

The test error was approximately

\[
\boxed{11.22\%}.
\]

The corresponding accuracy was approximately

\[
\boxed{88.78\%}.
\]

\subsection{(h) K-Nearest Neighbors}

KNN was evaluated using several values of $K$. Because KNN is based on
distances, the predictors were standardized before fitting the classifier.

\begin{lstlisting}
from sklearn.neighbors import KNeighborsClassifier
from sklearn.pipeline import make_pipeline
from sklearn.preprocessing import StandardScaler

for K in [1, 3, 5, 7, 9, 11]:

    knn = make_pipeline(
        StandardScaler(),
        KNeighborsClassifier(n_neighbors=K)
    )

    knn.fit(X_train, y_train)

    knn_pred = knn.predict(X_test)

    error = 1 - accuracy_score(y_test, knn_pred)

    print("K =", K, "Test error =", error)
\end{lstlisting}

The observed test errors were:

\begin{table}[H]
\centering
\begin{tabular}{cc}
\toprule
$K$ & Test Error \\
\midrule
1 & 7.65\% \\
3 & 8.16\% \\
5 & 9.18\% \\
7 & 10.20\% \\
9 & 11.22\% \\
11 & 11.22\% \\
\bottomrule
\end{tabular}
\caption{KNN test errors for the Auto dataset.}
\end{table}

The best performance was obtained using

\[
\boxed{K=1}
\]

with a test error of approximately

\[
\boxed{7.65\%}.
\]

\subsection{Exercise 14 Summary}

\begin{table}[H]
\centering
\begin{tabular}{lcc}
\toprule
Method & Test Error & Test Accuracy \\
\midrule
KNN ($K=1$) & 7.65\% & 92.35\% \\
LDA & 9.69\% & 90.31\% \\
Logistic Regression & 10.20\% & 89.80\% \\
Naive Bayes & 11.22\% & 88.78\% \\
QDA & 11.73\% & 88.27\% \\
\bottomrule
\end{tabular}
\caption{Comparison of classification methods for the Auto dataset.}
\end{table}

Among the methods tested, KNN with $K=1$ produced the lowest test error.
LDA was the second-best method, while QDA produced the highest test error
among the five approaches.

The results also illustrate an important point from classification theory:
greater model flexibility does not automatically result in better test
performance.

% =================================================
% EXERCISE 16
% =================================================

\section{Exercise 16: Boston Dataset}

The objective of this exercise is to predict whether a Boston suburb has a
crime rate above or below the median using classification methods.

\subsection{Creating the Binary Response}

The first step is to create a binary response variable called
\texttt{crim01}.

\begin{lstlisting}
from ISLP import load_data

Boston = load_data('Boston')

median_crim = Boston['crim'].median()

Boston['crim01'] = (
    Boston['crim'] > median_crim
).astype(int)
\end{lstlisting}

The response variable is interpreted as follows:

\begin{itemize}
    \item \texttt{crim01 = 0}: crime rate below or equal to the median.
    \item \texttt{crim01 = 1}: crime rate above the median.
\end{itemize}

The median crime rate was approximately

\[
\boxed{0.25651}.
\]

\subsection{Predictor Exploration}

The available predictors were examined to identify variables that were most
strongly associated with \texttt{crim01}.

The variables selected for the classification models were:

\[
\boxed{
\text{nox, rad, age, tax, indus, dis}
}
\]

These variables showed relatively strong associations with whether a suburb
had a crime rate above or below the median.

\subsection{Training and Test Sets}

The observations were divided into training and test sets using a 50--50
split, with \texttt{random\_state=1} and stratification.

\begin{lstlisting}
from sklearn.model_selection import train_test_split

features = [
    'nox',
    'rad',
    'age',
    'tax',
    'indus',
    'dis'
]

X = Boston[features]
y = Boston['crim01']

X_train, X_test, y_train, y_test = train_test_split(
    X,
    y,
    test_size=0.5,
    random_state=1,
    stratify=y
)
\end{lstlisting}

\subsection{Logistic Regression}

The logistic regression model produced a test error of approximately

\[
\boxed{15.42\%}.
\]

The corresponding test accuracy was approximately

\[
\boxed{84.58\%}.
\]

\subsection{Linear Discriminant Analysis}

LDA produced a test error of approximately

\[
\boxed{16.21\%}.
\]

The corresponding accuracy was approximately

\[
\boxed{83.79\%}.
\]

\subsection{Naive Bayes}

Gaussian naive Bayes produced a test error of approximately

\[
\boxed{17.00\%}.
\]

The corresponding accuracy was approximately

\[
\boxed{83.00\%}.
\]

\subsection{K-Nearest Neighbors}

Several values of $K$ were evaluated.

\begin{table}[H]
\centering
\begin{tabular}{cc}
\toprule
$K$ & Test Error \\
\midrule
1 & 6.72\% \\
3 & 7.91\% \\
5 & 7.51\% \\
10 & 8.30\% \\
\bottomrule
\end{tabular}
\caption{KNN test errors for the Boston dataset.}
\end{table}

The best result was obtained using

\[
\boxed{K=1}
\]

with a test error of approximately

\[
\boxed{6.72\%}.
\]

\subsection{Exercise 16 Summary}

\begin{table}[H]
\centering
\begin{tabular}{lcc}
\toprule
Method & Test Error & Test Accuracy \\
\midrule
KNN ($K=1$) & 6.72\% & 93.28\% \\
Logistic Regression & 15.42\% & 84.58\% \\
LDA & 16.21\% & 83.79\% \\
Naive Bayes & 17.00\% & 83.00\% \\
\bottomrule
\end{tabular}
\caption{Comparison of classification methods for the Boston dataset.}
\end{table}

KNN with $K=1$ achieved the lowest test error among the methods considered.
Its performance was substantially better than logistic regression, LDA, and
naive Bayes for this particular train-test split.

% =================================================
% OVERALL DISCUSSION
% =================================================

\section{Overall Discussion}

The four exercises demonstrate several important principles of
classification.

First, LDA and QDA differ mainly in their assumptions about the covariance
structure of the predictors. LDA assumes a common covariance matrix across
classes and therefore produces a linear decision boundary. QDA allows
class-specific covariance matrices and can therefore produce a quadratic
decision boundary.

Second, the additional flexibility of a model does not necessarily result in
better test performance. QDA can fit more complex relationships than LDA,
but this flexibility can increase variance and lead to poorer predictions
when the additional complexity is unnecessary.

Third, the computational exercises demonstrate that classification
performance depends strongly on the dataset. There is no single
classification method that is guaranteed to outperform all others in every
problem.

For the Auto dataset, KNN with $K=1$ achieved the lowest test error among the
methods evaluated. For the Boston dataset, KNN with $K=1$ also achieved the
lowest test error among the tested models.

The KNN results should nevertheless be interpreted carefully. KNN is a
non-parametric method and can have high variance, particularly when a small
value of $K$ is used. The ISLP text notes that the choice of $K$ has an
important effect on the flexibility of KNN and that $K=1$ can produce a very
flexible classifier.

% =================================================
% CONCLUSION
% =================================================

\section{Conclusion}

This chapter provided practical experience with several fundamental
classification methods: logistic regression, LDA, QDA, naive Bayes, and
KNN.

Exercise 5 demonstrated the bias--variance trade-off between LDA and QDA.
LDA is generally preferable when the true decision boundary is linear and
the sample size is relatively small, while QDA becomes more attractive when
the decision boundary is non-linear or when a sufficiently large sample is
available.

Exercise 10 showed mathematically that when there is one predictor, the LDA
log posterior odds are a linear function of that predictor.

Exercises 14 and 16 demonstrated the application of classification methods
to real datasets. In the particular train-test splits used in this analysis,
KNN with $K=1$ produced the lowest test error for both the Auto and Boston
datasets.

The exact test error rates depend on the training-test split. Therefore, the
reported values should be interpreted as the results for the specified
50--50 split with \texttt{random\_state=1}, rather than as universal
properties of the datasets.

Overall, the exercises demonstrate why model selection should be based on
test performance and an understanding of the assumptions and flexibility of
each classification method rather than simply choosing the most complex
model.

% =================================================
% REFERENCES
% =================================================

\section{References}

\begin{enumerate}

\item James, G., Witten, D., Hastie, T., Tibshirani, R., \& Taylor, J.
\textit{An Introduction to Statistical Learning with Applications in Python}.
Springer, 2023.

\item ISLP Python package and datasets used for the computational exercises.

\item Scikit-learn documentation for classification methods including
Logistic Regression, Linear Discriminant Analysis, Quadratic Discriminant
Analysis, Gaussian Naive Bayes, and K-Nearest Neighbors.

\end{enumerate}

\end{document}